% Profile: mistral-rse -- Research Software Engineer, Mistral AI (Paris /
% Research Software Engineering team, Paris + London).
%
% What the posting asks for, and what this profile puts in front of it:
%
%   "accelerate researchers by owning the complex parts of large-scale
%    pipelines and delivering robust internal tooling"
%        -> the summary's first sentence, and the CUDA/bindings and coding-agent
%           bullets in section_experience
%   "enforcing code review, CI, and observability standards"
%        -> the GitLab CI/CD bullet
%   "Python fluency plus proficiency in one systems language"
%        -> the reordered technology tag on the Fraunhofer entry
%   "GPU cluster or CUDA expertise" (nice-to-have)
%        -> CUDA kernels + DGX A100/H200 administration
%   "open-source contributions" (nice-to-have)
%        -> the summary's second sentence, which names what is public; the
%           upstreamed-fixes bullet in section_experience; and the header's
%           github.com/Kataglyphis, which is the only link a reader needs
%   Master's in Computer Science
%        -> section_education, which is why it stays in over Projects/Skills
%
% Not claimed: Kubernetes, SLURM. See the note in
% section_headline_mistral_rse.tex.
%
% Paris: section_private carries French at B1, which is the reason that section
% earns its place here rather than Projects. Do not drop it from this profile.

% The tagline has room for three things and CI/CD is not one of them: it is
% already a bullet under Fraunhofer IPA, and the header should not spend a
% third of its width restating the timeline. The slot goes to the AI/ML side
% of the job instead, which is both the part this employer cares about and the
% part the old tagline said nothing about. Do not put CI/CD back; the posting's
% CI requirement is answered by the bullet, not here.
%
% "LLM/VLM serving", not just "LLM serving": the vision-language half is its
% own claim and its own bullet (vLLM-served VLMs for auto-labeling, plus LoRA
% fine-tuning), and it is the half that matters most to an employer shipping
% open-weight vision models. The abbreviation earns its place -- spelled out,
% "Vision-Language-Modelle" pushes the German tagline past the header width.
%
% "GPU acceleration" and not "CUDA kernels": the broader term covers the whole
% span the bullets show -- custom kernels, C++ bindings, edge accelerators, the
% DGX nodes -- where "CUDA kernels" names only the narrowest of them. It briefly
% was "CUDA kernels", purely because the German would not fit otherwise; that is
% a typesetting reason, and typesetting reasons do not get to pick the wording.
%
% The German pays for "GPU-Beschleunigung" at the other end instead, with
% "Tooling" where the English says "research tooling". It has to pay somewhere:
% "Wissenschaftlicher Mitarbeiter @ Fraunhofer IPA | " eats twelve characters
% more of the header than its English counterpart, leaving about 42 for the
% three items. Past roughly 92 characters the line breaks mid-word
% ("Forschungstoo-") and STRICT_WARNINGS fails on the underfull hbox that
% follows -- which the old "GPU-Beschleunigung, CI/CD & Forschungstooling" was
% already doing, unnoticed, because nobody had rendered the German variant.
% Measure the whole line, not the part being edited, before growing either.
\newcommand{\cvTagline}{\IfLanguageName{english}{Research Associate @ Fraunhofer IPA | GPU acceleration, LLM/VLM serving \& research tooling}{Wissenschaftlicher Mitarbeiter @ Fraunhofer IPA | GPU-Beschleunigung, LLM/VLM \& Tooling}}

% The Gymnasium entry goes: this role requires a Master's, and those three
% lines buy the summary at the top of the page. See cv.tex.
\cvshowschoolfalse
% ... and both thesis lines. Together with the heading below, this is what pays
% for the tutor entry. Degrees, dates, grades and subject tags all stay.
\cvshowthesislinesfalse
% The tutor entry follows the other jobs under Experiences instead of opening a
% section of its own: it is a paid research-assistant post, it reads correctly
% there, and the heading it no longer prints is most of what it cost to fit.
\cvshowteachingtitlefalse
% The role is in Paris and this CV is in Stuttgart's letterhead; saying so in
% the header removes a question the reader would otherwise have to hold open.
% Deliberately not in the summary: that paragraph is about the work, and ending
% it on logistics would trade its last line for an administrative note.
\cvrelocatetrue

% Layout, not content: this profile carries one section body more than the
% page nominally fits, and the Languages/Hobbies/Publications row missed by
% about two lines. \emptySeparator is a full blank table row between
% entries; pulling half a line back out of each recovers that without
% dropping anything. The README sanctions exactly this -- a profile may
% redefine what the class exposes, for layout, not for content. If an edit
% makes this profile spill again, take the line back here first, before
% touching a bullet.
\renewcommand\emptySeparator{\multicolumn{2}{c}{}\\[-3\baselineskip]}

\newcommand{\cvBody}{%
    % Summary for the mistral-rse profile: Research Software Engineer, Mistral AI
% (Paris). Written against that posting's first responsibility -- "accelerate
% researchers by owning the complex parts of large-scale pipelines and
% delivering robust internal tooling" -- and says nothing the experience and
% education sections below do not back up.
%
% Deliberately absent: Kubernetes and SLURM. The posting asks for them and the
% claim would not be true; the reproducible-container and GPU-node work is what
% is offered against that requirement instead. Do not "improve" this by adding
% an orchestrator.
%
% Three verbs, not one. The middle of this paragraph once read "I own the CUDA
% kernels ... , the institute's CI/CD, its DGX A100/H200 nodes and its
% vision-language models" -- one verb stretched over four unlike things, and it
% fitted none of them well: you write kernels, you run nodes, you serve models.
% Write/run/serve says what the work actually is at no extra length.
%
% The paragraph is about Fraunhofer IPA, deliberately. It used to spend its
% whole second half on personal repositories -- "my renderers, my ML code and
% this page's pipeline are public" -- which is roughly half the space a reader
% gives this CV, handed to spare-time work. The open-source claim did not have
% to go with them: the fixes to the institute's toolchain are *the job*, not a
% hobby, so saying it that way keeps the claim, moves it onto paid work, and
% buys the line the vision-language clause needed. github.com/Kataglyphis in
% the header, the Hobbies column and the footer's typeset-with credit carry
% what is left of the personal side; they do not need a sentence here too.
%
% "not into a local patch" is load-bearing and has been lost to a length trim
% once already. The claim is a habit with a consequence, not "passionate about
% open source" -- the contrast is what makes it a practice rather than a taste.
% Backed by the Experiences bullet on upstreamed fixes, which no longer names
% the two PRs it used to, on request. Worth its line for this employer in
% particular: Mistral ships open-weight models, and the Fraunhofer bullets are
% careful to say open-weight where that is what they mean.
%
% Every item here is named, not explained: each one is a bullet in
% section_experience immediately below, and a summary that re-explains the page
% under it is the first thing a reader skips. The one clause that goes further
% -- "that replaced its hand-labeling" -- is there because it is an outcome,
% not a restatement, and because it is the only thing in the paragraph that
% says what "faster" bought anyone.
%
% One page is a hard requirement and this paragraph is the first thing to break
% it: three lines fit, four push the Languages/Hobbies/Publications row onto
% page two -- which has already happened here twice. The ceiling is about 345
% rendered characters in either language, and German reaches it sooner, which
% is why it says "die das Labeling übernehmen" where the English can afford
% "that replaced its hand-labeling". Rebuild BOTH languages after any edit.
% No GitHub URL here: the header already carries it.
\IfLanguageName{english}{
    \par{I make other people's research run faster: at Fraunhofer IPA I write the CUDA kernels and Python bindings behind the team's methods, run its CI/CD and DGX A100/H200 nodes, and serve the vision-language models that replaced its hand-labeling. Where the tooling we depend on is open source, my fixes go upstream, not into a local patch.}
}{
    \par{Ich mache die Forschung anderer schneller: Am Fraunhofer IPA schreibe ich CUDA-Kernels und Python-Bindings hinter den Verfahren des Teams, betreibe CI/CD und DGX-A100/H200-Knoten und hoste die Vision-Language-Modelle, die das Labeling übernehmen. Wo unsere Werkzeuge Open Source sind, gehen meine Fixes upstream statt in lokale Patches.}
}
%
    \vspace{0.25cm}%
    % YAAC Another Awesome CV LaTeX Template
%
% This template has been downloaded from:
% https://github.com/darwiin/yaac-another-awesome-cv
%
% Author:
% Christophe Roger
%
% Template license:
% CC BY-SA 4.0 (https://creativecommons.org/licenses/by-sa/4.0/)

%Section: Scholarships and additional info
\IfLanguageName{english}{
    
    \sectionTitle{Experiences}{\faHistory}
        \begin{experiences}
        
        \experience
        {\DTMlangsetup*{showdayofmonth=false}Present}{Research Associate | Scene Analysis \& Activity Recognition}{Stuttgart}{Fraunhofer IPA}
        {\DTMlangsetup*{showdayofmonth=false}\DTMdisplaydate{2024}{02}{01}}
        {   \begin{itemize}
                % ORDER IS PER PROFILE (since 2026-09-22). Each bullet below is
                % defined once, in place, as \cvFhg<Name>, with its notes beside
                % it, and nothing prints until \cvFraunhoferOrder at the end of
                % this list. cv.tex sets the default order, which is the order
                % written here. A profile may \renewcommand it to lead with what
                % its posting asks for first, and terabase-cv does. Reordering
                % selects and orders; it does not fork the prose. Keep the
                % German branch's \def names identical.
                %
                % Ordered for this reader, not chronologically: the platform
                % cluster first -- CI/CD, builds, gateway, docs -- then the
                % research work. Re-ordered when \ifcvshowcudakernels went in:
                % hiding the kernels bullet on a platform profile silently
                % promoted the labeling result to the opening line, which is not
                % what a platform reader is looking for. The research profiles
                % still show the kernels bullet, so it still leads there.
                %
                % Was eight bullets; a reader gives an entry four or five before
                % skimming, and the last three were carrying real work into that
                % blind spot. Two merges did it, and both removed a repetition
                % rather than a fact:
                %   - the camera-pipeline bullet and the "computer vision
                %     workloads across GPU servers and edge devices" bullet were
                %     the same work seen twice, once from the client end and once
                %     from the hardware end. Joined, they show the whole span --
                %     which is the point, and which neither said alone.
                %   - CI/CD and DGX administration are one sentence: the platform
                %     the group runs on. The old CI/CD bullet also spent half a
                %     line explaining that CI/CD means build, test and deploy.
                %
                % "Replaced ... hand-labeling" and not "Automated data labeling":
                % same fact, but one names what it cost the group before. Verbs
                % that state an outcome are the only thing standing in for the
                % numbers this list still does not have.
                % The GLM-5.2 optimisations were a parent bullet plus two
                % sub-bullets: six lines for what fits in three as one sentence.
                % A nested itemize also indents twice, so each child line is
                % narrower and wraps sooner -- it costs lines twice over. Same
                % numbers, same order, no sub-list.
                \def\cvFhgKernels{\ifcvshowcudakernels
                \item Accelerated the team's efficient-learning and optimization methods with custom CUDA kernels and Python bindings to C++
                \fi}
                % Fraunhofer IPA's quantum-computing group. Added for the amd-hpc
                % profile (see the flag in cv.tex); ~700x is the owner's measured
                % figure. The link is the public project page, which is what makes
                % the claim checkable -- keep it. Replaces the RemaNet bullet on
                % that profile: edge-AI product work is the one item here an HPC
                % posting does not buy.
                \def\cvFhgQuantum{\ifcvshowquantum
                \item Ported a \textbf{quantum-physics simulation} from CPU to GPU for the \link{https://www.kqcbw.de}{KQCBW} with \textbf{CUDA} and \textbf{PyTorch}, profiling with \textbf{Nsight} and \textbf{VTune} to a \textbf{$\sim$700$\times$} speedup
                \fi}
                % Opens the platform cluster -- and on a platform profile, where
                % the kernels bullet above is hidden, opens the entry. It was the
                % only bullet here carrying no number, which for a platform
                % posting was the wrong one to leave bare. Ratio deliberately
                % unpinned: AGENTS.md forbids quoting a wall-clock digit that no
                % gate re-measures. Cache hit rate is measurable per run and is
                % the figure to swap in if one is ever wanted.
                % \ifcvshowmldetail clause (owner, 2026-09-22): the group
                % tracked experiments in MLflow and ClearML -- the posting's
                % "experiment tracking". Passive on purpose: these were the
                % team's tools, used, not owned like the CI/CD before them.
                \def\cvFhgCicd{%
                \item Owned the institute's GitLab CI/CD, DGX A100/H200 GPUs and \textbf{Grafana/Prometheus} monitoring; \textbf{BuildKit} and \textbf{sccache} caching cut builds from hours to minutes\ifcvshowmldetail; experiments tracked in \textbf{MLflow} and \textbf{ClearML}\fi
                }
                % Naming the projects is the point: "fixes go upstream" is a
                % posture, named projects are checkable evidence of root-cause
                % work in somebody else's codebase -- the posting's "root cause
                % analysis" ask, answered nowhere else on the page.
                %
                % LLVM and sccache, and "upstreamed" rather than a PR count:
                % owner decision 2026-09-04. ANTfrastructure's
                % docs/upstream-windows-patches.md "Already filed" table (checked
                % 2026-09-02) records LLVM #219275/#219276/#219200 and hcsshim
                % #2855 as OPEN, and lists sccache #2808 as an issue we opened
                % rather than a patch of ours -- but that register tracks only
                % what this repo filed, and it is the owner's own contribution
                % history, not a three-day-old table, that settles this.
                %
                % Two things follow. "Upstreamed" reads as MERGED, so it has to
                % be merged if a reviewer opens the PR list. And do not "restore"
                % hcsshim or a PR count here from the register: both were
                % deliberately removed.
                % Behind \ifcvshowbuilds (default on; see cv.tex).
                \def\cvFhgBuilds{\ifcvshowbuilds
                \item Advanced \textbf{CMake} build systems for C++ codebases, \textbf{GStreamer} (\textbf{Meson}) and \textbf{LiteRT-LM} (\textbf{Bazel}) built from source; fixes upstreamed to \textbf{LLVM} and \textbf{sccache}, not kept as local patches
                \fi}
                % Third in the platform cluster, and the only place the CV claims
                % Go -- see the flag in cv.tex. "45 researchers" is the scale
                % figure: the plugins are infrastructure a whole department runs
                % through, which a bare "token metering" does not convey.
                \def\cvFhgGateway{\ifcvshowgateway
                \item Extended the institute's \textbf{AI gateway} with custom \textbf{WebAssembly} plugins in \textbf{Go}: token metering and rate limiting for its \textbf{45 researchers}
                \fi}
                % Closes the platform cluster, deliberately after the systems it
                % documents: the posting asks for documentation "that will help
                % them make the most out of the tools and systems you'll build",
                % which is this order exactly. It names those systems rather than
                % saying "docs" -- an abstract documentation claim is the easiest
                % thing on a CV for a reader to discount.
                \def\cvFhgDocs{\ifcvshowdocs
                \item Wrote the runbooks and onboarding docs for the CI/CD and GPU cluster the group works from
                \fi}
                % "Replaced ... hand-labeling" and not "Automated data labeling":
                % same fact, but one names what it cost the group before. Verbs
                % that state an outcome are the only thing standing in for the
                % numbers this list still does not have.
                % \ifcvshowmldetail clause (owner, 2026-09-22): the pipelines
                % run in the owner's own research tool for self-labeling, and
                % its labels partly became the training data of the owner's
                % projects -- "data annotation tooling" and dataset curation.
                % Not claimed: bounding boxes or a human review step (not
                % asked, not answered).
                % Reworded 2026-09-23 by the owner: the point is that some of
                % the owner's projects used the generated labels, not that
                % "its labels partly trained my models". The first clause,
                % kept:
                % \ifcvshowmldetail{} in my own research tool for self-labeling; its labels partly trained my models\fi
                \def\cvFhgVlm{%
                \item Cut hand-labeling time by \textbf{75\%} with \textbf{vLLM-served vision-language model pipelines}\ifcvshowmldetail{} in my own research tool for self-labeling; some of my projects used its generated labels\fi
                }
                % Both GLM-5.2 bullets are the same box: 6 GPUs of one H200
                % system. It said "6x H100" here until now, which read as a
                % second, different deployment. The hardware is named once, in
                % this bullet; the one below says "there" rather than repeating
                % it. Keep it that way -- naming it twice invites the reader to
                % assume two clusters, which is what the mismatch already did.
                % One deployment, three optimisations -- not two bullets. It was
                % split as "Deployed vLLM serving for GLM-5.2 on 6 GPUs..." and
                % "Ran GLM-5.2 (744B MoE) there at NVFP4...", which named the
                % model twice, the machine twice, and used two verbs for one
                % deploy. Six lines for what fits in four, and the split invited
                % the reader to count two systems. Every number is still here.
                % One headline number per optimisation -- and since 2026-09-04
                % only that: "-46% latency" came out because the bullet had
                % drifted back to four figures, and 10x throughput already makes
                % the speed claim. Two numbers a reader carries away beat four
                % they skim. The supporting figures were dropped earlier for the
                % same reason, not because they are weak:
                % ~143 tok/s at 55k context, 65% draft acceptance, -60% worst-
                % case latency, "no loss on large prefills", 8/8 perf gates, and
                % "than BF16" on the memory figure. Eight numbers in one bullet
                % is a spec sheet; three is a claim a reader carries away, and
                % the rest is what an interview is for. All of it is one edit
                % away if the page ever grows.
                % The coding agent IS this deployment -- it runs on the GLM-5.2
                % instance described here, so listing it separately claimed two
                % pieces of work where there is one, and the reader had no way
                % to tell. Purpose first, engineering after: what the cluster is
                % for, then what was done to make it fast enough to be usable.
                % "Started ... on our own initiative ... and grew it into
                % production" carries three things a claim never could: the work
                % was not assigned, it was done with someone, and it was carried
                % all the way. The posting asks for "collaborative, low-ego" --
                % this is the evidence, and evidence is the only way to answer
                % that ask without contradicting it by asserting it.
                % The opening stays -- unassigned, with someone, carried to
                % production is what this bullet is really evidence for. The
                % technical tail was a spec sheet: three techniques, each with
                % its own parenthesised metrics, five numbers in all. Kept the
                % two custom pieces of vLLM work and the memory result; dropped
                % TP2xPP3, "median", and the throughput/latency split per
                % technique. All of it is in the git history if it should return.
                % "a production service the institute's researchers use" and not
                % "production": the platform posting asks for environments built
                % FOR researchers, and this bullet is the only place the CV can
                % name its users. Says nothing new -- an institute coding agent
                % in production has exactly those users. An adoption figure
                % would be stronger still; there is no measured one to cite.
                % Before "privacy-first" (2026-09-22), kept for reference:
                % \item Started the institute's coding agent with a colleague on our own initiative and grew it into a production service its researchers use: self-hosted \textbf{GLM-5.2 and GLM-5.3} on \textbf{6 H200 GPUs}, with custom \textbf{speculative decoding} and \textbf{scheduling} (\textbf{$\sim$10$\times$ throughput}) and \textbf{NVFP4 + FP8 KV cache} (\textbf{1M-token context})
                % "privacy-first": the owner's word (2026-09-22) for why it is
                % self-hosted -- AI-first engineering without sending code to a
                % third party. It fits in the last line's slack.
                % "GLM-5.3 with vision capabilities" instead of "GLM-5.2 and
                % GLM-5.3": the owner's wording, 2026-09-23. It names the model
                % the service runs now and says it sees as well as writes,
                % which ties the agent to the vision work above it. GLM-5.2
                % stays true history (see the notes above). The version with
                % both models, kept:
                % \item Started the institute's coding agent with a colleague on our own initiative and grew it into a privacy-first production service its researchers use: self-hosted \textbf{GLM-5.2 and GLM-5.3} on \textbf{6 H200 GPUs}, with custom \textbf{speculative decoding} and \textbf{scheduling} (\textbf{$\sim$10$\times$ throughput}) and \textbf{NVFP4 + FP8 KV cache} (\textbf{1M-token context})
                \def\cvFhgAgent{%
                \item Started the institute's coding agent with a colleague on our own initiative and grew it into a privacy-first production service its researchers use: self-hosted \textbf{GLM-5.3 with vision capabilities} on \textbf{6 H200 GPUs}, with custom \textbf{speculative decoding} and \textbf{scheduling} (\textbf{$\sim$10$\times$ throughput}) and \textbf{NVFP4 + FP8 KV cache} (\textbf{1M-token context})
                }
                % RemaNet = Remanufacturing Network, Horizon Europe grant
                % 101138627, 25 partners, Jan 2024 - Dec 2026, Fraunhofer IPA a
                % partner -- which is why it lines up with the February 2024
                % start date. Named because it is checkable: a reader can look
                % the project up, which "AI camera pipelines" alone does not
                % allow. The grant number stays out; the programme is the signal.
                % "Raspberry Pi/Hailo edge devices" undersold this: the Hailo
                % went into a Berghof PLC, i.e. an industrial controller was
                % retrofitted with an AI accelerator. That is a fact a reader can
                % picture and few CVs can claim; "edge devices" is a category.
                % Paid for by dropping "low-latency, high-throughput" and "end to
                % end" -- self-assessed adjectives, and the list from GPU server
                % down to PLC and up to the clients already says "end to end".
                % The mistral-rse builds have ~0pt spare, so this had to come out
                % no longer than it went in; check both before growing it.
                % The frontend is RemaNet's, not the PLC's. It used to trail the
                % sentence -- "...a Berghof PLC I fitted with a Hailo
                % accelerator, with a React web frontend" -- where it attaches
                % to the nearest noun, which is the PLC, and the second "with"
                % made that reading worse. Naming it first and hanging the
                % inference under it reads as one stack, top down, and everything
                % in the list belongs to RemaNet because RemaNet precedes the colon.
                % Linked like KQCBW in the quantum bullet: both are research
                % consortia a reader can look up, so both get the link colour
                % rather than bold. remanet-project.eu is the official site; it
                % names grant 101138627 and lists Fraunhofer IPA as a partner
                % (checked 2026-09-22). Not bold: the colour already marks the
                % name, and bold plus colour would outweigh KQCBW.
                % The bold, unlinked version, kept for reference:
                % \item Shipped \textbf{AI camera pipelines} for \textbf{RemaNet} (Horizon Europe): a \textbf{React} web frontend over C++ inference on NVIDIA GPU servers and a \textbf{Berghof PLC I fitted with a Hailo accelerator}
                %
                % CORRECTED 2026-09-22, by the owner. The Berghof PLC was NOT part
                % of RemaNet; it was a separate project. So was a third one, an
                % industry contract. The owner ran the whole ML lifecycle (design,
                % training, evaluation, deployment) in all three:
                %   RemaNet  -- activity detection; React frontend over C++
                %               inference on NVIDIA GPU servers (owner, later
                %               the same day: ONNX inference, PyTorch training,
                %               NVIDIA GPUs)
                %   Berghof  -- object detection on an embedded PLC fitted with a
                %               Hailo AI accelerator; the embedded target with a
                %               special AI chip was the challenge
                %   industry -- person detection for safety: 4 cameras at once,
                %               their 4 inferences batched
                % The notes above ("everything in the list belongs to RemaNet
                % because RemaNet precedes the colon") describe the superseded
                % bullet below. Its wording was right; the attribution was wrong.
                % Each project now has its own bullet and its own flag:
                % \ifcvshowremanet, \ifcvshowberghof and \ifcvshowsafety (see
                % cv.tex).
                % Superseded 2026-09-22 (it attributes the Berghof PLC to RemaNet):
                % \item Shipped \textbf{AI camera pipelines} for \link{https://remanet-project.eu}{RemaNet} (Horizon Europe): a \textbf{React} web frontend over C++ inference on NVIDIA GPU servers and a \textbf{Berghof PLC I fitted with a Hailo accelerator}
                % Before "25 partners" (2026-09-22), kept for reference:
                % \item Shipped \textbf{activity detection} for \link{https://remanet-project.eu}{RemaNet} (Horizon Europe): camera pipelines with a \textbf{React} web frontend over C++ inference on NVIDIA GPU servers
                % "25 partners": CORDIS and remanet-project.eu both list 25
                % organisations from 8 countries (checked 2026-09-22). It is the
                % cheapest evidence on the page of delivering inside a
                % multi-organisation consortium, i.e. working with external
                % partners. It stays inside the second line's slack, so it adds
                % no line to any profile.
                % With "25 partners" but before the owner's reformulation
                % (2026-09-22), kept for reference:
                % \item Shipped \textbf{activity detection} for \link{https://remanet-project.eu}{RemaNet} (Horizon Europe, 25 partners): camera pipelines with a \textbf{React} web frontend over C++ inference on NVIDIA GPU servers
                % Reformulated by the owner, 2026-09-22: inference runs on ONNX
                % and training in PyTorch, both on NVIDIA GPUs. "C++ inference"
                % is not wrong, but it named the language where the owner names
                % the runtime and the training framework. C++ stays in the tag
                % row. \slash, not "/": a bare slash glues "inference/PyTorch"
                % into one unbreakable word, which TeX can only fix with a bad
                % line. The owner's text ended in ";", which is dropped: no
                % other bullet ends in punctuation.
                \def\cvFhgRemanet{\ifcvshowremanet
                \item Shipped \textbf{activity detection} for \link{https://remanet-project.eu}{RemaNet} (Horizon Europe, 25 partners): camera pipelines with a \textbf{React} web frontend and \textbf{ONNX} inference\slash\textbf{PyTorch} training on NVIDIA GPUs
                \fi}
                % Before the frame rate and the framework (2026-09-22), kept:
                % \item Ran \textbf{object detection} on an embedded \textbf{Berghof PLC I fitted with a Hailo AI accelerator}
                % "~14 fps" and "PyTorch-trained": both from the owner,
                % 2026-09-22. The frame rate is measured on the Hailo chip. The
                % models were trained in PyTorch and deployed to the embedded
                % target -- the embedded chip was the challenge of this
                % project, and the pair says so in two words. Two lines now,
                % not one; the profiles that show this bullet have the room
                % (mistral-rse and amd-hpc hide it).
                % \ifcvshowmldetail clause (owner, 2026-09-22): the models ran
                % INT8-quantized on the Hailo chip, and the owner iterated with
                % the industry customers on field feedback -- the posting's
                % "gathering field feedback ... closing the loop between model
                % performance and real-world outcomes", answered with the
                % project where it happened most.
                % The clause's first wording, kept: it set "feedback" alone on a
                % third line, one line for one word.
                % \ifcvshowmldetail; \textbf{INT8}-quantized and iterated with industry customers on field feedback\fi
                \def\cvFhgBerghof{\ifcvshowberghof
                \item Ran PyTorch-trained \textbf{object detection} at \textbf{$\sim$14 fps} on an embedded \textbf{Berghof PLC I fitted with a Hailo AI accelerator}\ifcvshowmldetail; \textbf{INT8}-quantized, iterated on customer field feedback\fi
                \fi}
                % Earlier wordings, 2026-09-22, kept for reference:
                % \item Built \textbf{person detection} for industrial safety: \textbf{4 cameras} in parallel, their inferences \textbf{batched}
                % \item Built \textbf{person detection} for industrial safety: \textbf{4 cameras} in parallel, inferences \textbf{batched}
                % "Delivered ... for an industrial customer": it was a custom
                % industry project, so the model went to a paying customer and
                % had to hold up there. That is the "building, shipping and
                % supporting production B2B applications" a posting may ask for,
                % and "Built" did not say it.
                \def\cvFhgSafety{\ifcvshowsafety
                \item Delivered \textbf{person detection} for an industrial customer's safety application: \textbf{4 cameras} in parallel, inferences \textbf{batched}
                \fi}
                % KI-PrOT (owner, 2026-09-22; ki-prot.de checked the same day):
                % "Optimierte Produktionsplanung für die Oberflächentechnik",
                % funded by the BMFTR, coordinated by Softec, with Fraunhofer
                % IPA developing the AI methods. The owner built the visual
                % classification of parts into drag-out classes
                % (Verschleppungsklassen), created and annotated its dataset,
                % and versioned it with DVC. The site's stated aim is to reduce
                % chemical consumption; that is a goal, not a measured result,
                % so the bullet says "to reduce". It is the CV's only
                % classification work and its only DVC. Behind \ifcvshowkiprot,
                % default off.
                \def\cvFhgKiprot{\ifcvshowkiprot
                \item Built \textbf{visual classification} of parts into drag-out classes for \link{https://ki-prot.de}{KI-PrOT} to reduce chemical use in surface treatment: dataset created, annotated and versioned with \textbf{DVC}
                \fi}
                % Everything above is only defined. This line prints it, in the
                % order cv.tex or the profile sets.
                \cvFraunhoferOrder
            \end{itemize}
        }
        % Tags follow the bullets above, strongest first for this reader.
        % Dropped: "Edge AI", a category where the bullet now names the actual
        % accelerator, and "DGX A100/H200", which the CI/CD bullet and the GLM
        % one already spell out -- a tag that restates the prose above it only
        % costs width. Added NVFP4/FP8 (the quantisation work had no tag at all),
        % React and Hailo. LoRA stays although the prose no longer says it: the
        % fine-tuning is real, and a tag is the cheapest place to keep a keyword.
        % One line is the budget -- this row wrapping costs the mistral-rse
        % builds their last few points of slack.
        % PyTorch added 2026-09-22: it is on this entry's page (the quantum
        % port) and several postings list it as a must-have, and the row
        % still fits on one line. The list before, kept for reference:
        % {Python, Go, C++, CUDA, CMake, Bazel, Meson, vLLM, LoRA, GitLab CI/CD}
        % uv added the same day (owner): Python package and environment
        % management with Astral's uv, used at Fraunhofer. It sits next to
        % Python on purpose -- it is Python tooling. The tag is the tool's
        % own name, "uv". "Astral uv" was tried first and wrapped
        % "GitLab CI/CD" onto a second tag line, which would have cost every
        % profile a line, mistral-rse its page. The lists before, kept:
        % {Python, Go, C++, CUDA, PyTorch, CMake, Bazel, Meson, vLLM, LoRA, GitLab CI/CD}
        % {Python, Astral uv, Go, C++, CUDA, PyTorch, CMake, Bazel, Meson, vLLM, LoRA, GitLab CI/CD}
        {Python, uv, Go, C++, CUDA, PyTorch, CMake, Bazel, Meson, vLLM, LoRA, GitLab CI/CD}
        
        \emptySeparator
        
        \experience
        {\DTMlangsetup*{showdayofmonth=false}\DTMdisplaydate{2023}{05}{01}}{Research assistant}{Karlsruhe}{CAS Software AG}
        {\DTMlangsetup*{showdayofmonth=false}\DTMdisplaydate{2022}{10}{15}}  
        {    \ifcvshowcasdetail
            \begin{itemize}
                % Was three lowercase sentence fragments, one of which
                % ("providing user guidance in mixed reality applications")
                % only restated the entry it sits in.
                % Was two bullets that both ended on "guidance": the ONNX one
                % and "Applied eye tracking and then-new ML approaches to the
                % guidance problem". Eye tracking was the only fact the second
                % one added, so it moved into the first. One line back in German,
                % where the pair wrapped; in English it reads tighter either way.
                % One line since 2026-09-22 (the terabase-cv owner wanted this
                % entry described again, and the page had room for one line,
                % not two). Same facts; "with eye tracking" also stops reading
                % as if the models had it. The two-line version, kept:
                % \item Trained and deployed ONNX models with eye tracking for real-time user guidance in mixed reality
                \item Trained and deployed ONNX models for real-time eye-tracking guidance in mixed reality
            \end{itemize}
            \fi
        }
        {Machine learning, ONNX, Mixed reality, HoloLens2, C\#, Python}
        
        \end{experiences}
    }{
    
    \sectionTitle{Berufserfahrung}{\faHistory}
    \begin{experiences}
    % german
    \experience
    {\DTMlangsetup*{showdayofmonth=false}heute}{Wissenschaftlicher Mitarbeiter | Szenenanalyse \& Aktivitätserkennung}{Stuttgart}{Fraunhofer IPA}
    {\DTMlangsetup*{showdayofmonth=false}\DTMdisplaydate{2024}{02}{01}}
    {   \begin{itemize}
            % Wording is tighter than the English throughout this list, not for
            % taste: German sets ~10% longer and this entry is what decides
            % whether the Sprachen/Hobbies/Publikationen row fits on page one.
            % Reihenfolge und Zusammenlegungen wie im englischen Zweig -- siehe
            % die Begruendung dort. Der Nominalstil bleibt: er ist im deutschen
            % Lebenslauf ueblich und hier ausserdem das, was in die Zeile passt.
            % Siehe die Anmerkung im englischen Zweig: keine Unterliste.
            % REIHENFOLGE JE PROFIL (seit 2026-09-22), siehe die Anmerkung im
            % englischen Zweig: jeder Punkt wird hier nur als \cvFhg<Name>
            % definiert, ausgegeben wird er erst von \cvFraunhoferOrder am Ende
            % der Liste. Die Namen müssen denen im englischen Zweig gleichen.
            \def\cvFhgKernels{\ifcvshowcudakernels
            \item \textbf{CUDA-Kernels} und \textbf{Python-Bindings} an C++, die die Lern- und Optimierungsverfahren des Teams beschleunigen
            \fi}
            % Siehe die Anmerkung im englischen Zweig: Quantencomputing-Gruppe
            % des Fraunhofer IPA, gemessener ~700x-Speedup, der Link ist die
            % öffentliche Projektseite. Ersetzt dort den RemaNet-Punkt.
            \def\cvFhgQuantum{\ifcvshowquantum
            \item \textbf{Quantenphysikalische Simulation} des \link{https://www.kqcbw.de}{KQCBW} von CPU auf GPU portiert (\textbf{CUDA}, \textbf{PyTorch}) und mit \textbf{Nsight} und \textbf{VTune} auf \textbf{$\sim$700$\times$} Speedup profiliert
            \fi}
            % OFFEN: eine Zeile bleibt gedehnt (badness 5667), STRICT_WARNINGS
            % schlägt fehl. \allowbreak am Schrägstrich ändert den Umbruch nicht;
            % "Monitoring mit Grafana/Prometheus" löst es, kostet aber die Seite.
            \def\cvFhgCicd{%
            \item Verantwortlich für GitLab-CI/CD, DGX-A100/H200-GPUs und \textbf{Grafana/Prometheus}-Monitoring; \textbf{BuildKit}- und \textbf{sccache}-Caching verkürzten Builds von Stunden auf Minuten\ifcvshowmldetail; Experiment-Tracking mit \textbf{MLflow} und \textbf{ClearML}\fi
            }
            % Siehe die Anmerkung im englischen Zweig.
            % Hinter \ifcvshowbuilds (standardmäßig an; siehe cv.tex).
            \def\cvFhgBuilds{\ifcvshowbuilds
            \item Fortgeschrittene \textbf{CMake}-Builds für C++-Codebasen, \textbf{GStreamer} (\textbf{Meson}) und \textbf{LiteRT-LM} (\textbf{Bazel}) aus dem Quellcode gebaut; Fixes an \textbf{LLVM} und \textbf{sccache} upstream eingebracht, statt sie lokal zu patchen
            \fi}
            % Siehe die Anmerkung im englischen Zweig.
            \def\cvFhgGateway{\ifcvshowgateway
            \item \textbf{KI-Gateway} des Instituts um eigene \textbf{WebAssembly}-Plugins in \textbf{Go} erweitert: Token-Metering und Rate-Limiting für \textbf{45 Forschende}
            \fi}
            % Siehe die Anmerkung im englischen Zweig: nennt die Systeme.
            \def\cvFhgDocs{\ifcvshowdocs
            \item Runbooks und Onboarding-Dokumentation für die CI/CD und den GPU-Cluster geschrieben, mit denen die Gruppe arbeitet
            \fi}
            % Siehe die Anmerkung im englischen Zweig. Umformuliert 2026-09-23;
            % die erste Fassung des Zusatzes, zum Nachschlagen erhalten:
            % \ifcvshowmldetail{} in meinem eigenen Forschungstool für Self-Labeling; seine Labels flossen teils ins Training meiner Modelle\fi
            \def\cvFhgVlm{%
            \item Manuelles Labeling um \textbf{75\%} reduziert: \textbf{vLLM-gehostete Vision-Language-Pipelines}\ifcvshowmldetail{} in meinem eigenen Forschungstool für Self-Labeling; einige meiner Projekte nutzten die generierten Labels\fi
            }
            % Siehe die Anmerkung im englischen Zweig: eine Maschine, einmal benannt.
            % Siehe die Anmerkung im englischen Zweig: ein Deployment, drei Optimierungen.
            % Siehe die Anmerkung im englischen Zweig: eine Kennzahl je Optimierung.
            % Siehe die Anmerkung im englischen Zweig: ein Deployment, nicht zwei.
            % Siehe die Anmerkung im englischen Zweig.
            % Siehe die Anmerkung im englischen Zweig.
            % Siehe die Anmerkung im englischen Zweig.
            % Vor "datenschutzkonform" (2026-09-22), zum Nachschlagen erhalten:
            % \item Coding-Agent des Instituts auf Eigeninitiative mit einem Kollegen als Dienst für die Forschenden in Produktion gebracht: selbst gehostetes \textbf{GLM-5.2 und GLM-5.3} auf \textbf{6 H200-GPUs}, mit \textbf{Speculative Decoding} und \textbf{Scheduling} (\textbf{$\sim$10$\times$ Durchsatz}) und \textbf{NVFP4 + FP8-KV-Cache} (\textbf{1M-Token-Kontext})
            % "GLM-5.3 mit Vision-Fähigkeiten" statt "GLM-5.2 und GLM-5.3"
            % (2026-09-23), siehe die Anmerkung im englischen Zweig. Die
            % Fassung mit beiden Modellen, zum Nachschlagen erhalten:
            % \item Coding-Agent des Instituts auf Eigeninitiative mit einem Kollegen als datenschutzkonformen Dienst für die Forschenden in Produktion gebracht: selbst gehostetes \textbf{GLM-5.2 und GLM-5.3} auf \textbf{6 H200-GPUs}, mit \textbf{Speculative Decoding} und \textbf{Scheduling} (\textbf{$\sim$10$\times$ Durchsatz}) und \textbf{NVFP4 + FP8-KV-Cache} (\textbf{1M-Token-Kontext})
            \def\cvFhgAgent{%
            \item Coding-Agent des Instituts auf Eigeninitiative mit einem Kollegen als datenschutzkonformen Dienst für die Forschenden in Produktion gebracht: selbst gehostetes \textbf{GLM-5.3 mit Vision-Fähigkeiten} auf \textbf{6 H200-GPUs}, mit \textbf{Speculative Decoding} und \textbf{Scheduling} (\textbf{$\sim$10$\times$ Durchsatz}) und \textbf{NVFP4 + FP8-KV-Cache} (\textbf{1M-Token-Kontext})
            }
            % Kein "end to end" hier: im Deutschen ein Anglizismus, den die
            % Doppelpunkt-Liste ohnehin schon sagt. Die Geraete stehen in einer
            % Klammer statt als Bindestrich-Kompositum: "Raspberry-Pi-/Hailo-
            % Edge-Geraeten" ist eine einzige unumbrechbare Einheit, und mit ihr
            % in der Zeile fand TeX keinen Umbruch mehr -- die erste Zeile wurde
            % auf Wortabstaende von mehreren Millimetern gezogen (Underfull
            % \hbox, badness 10000, was STRICT_WARNINGS zu Recht ablehnt).
            % NVIDIA faellt weg; der DGX-Punkt darunter nennt den Hersteller.
            % Siehe die Anmerkung im englischen Zweig.
            % "darauf" gestrichen: der Punkt lief um sechs Zeichen über zwei
            % Zeilen und kostete eine dritte fuer "native Apps" allein.
            % Siehe die Anmerkung im englischen Zweig.
            % Siehe die Anmerkung im englischen Zweig: das Frontend gehört zu RemaNet, nicht zur SPS.
            % Siehe die Anmerkung im englischen Zweig: verlinkt wie KQCBW.
            % Die fette, unverlinkte Fassung, zum Nachschlagen erhalten:
            % \item \textbf{KI-Kamera-Pipelines} für \textbf{RemaNet} (Horizon Europe): \textbf{React}-Web-Frontend über C++-Inferenz auf GPU-Servern und einer \textbf{Berghof-SPS mit Hailo}
            % KORRIGIERT 2026-09-22 (siehe die Anmerkung im englischen Zweig): die
            % Berghof-SPS war ein eigenes Projekt, nicht Teil von RemaNet.
            % Ersetzt 2026-09-22 (ordnet die Berghof-SPS RemaNet zu):
            % \item \textbf{KI-Kamera-Pipelines} für \link{https://remanet-project.eu}{RemaNet} (Horizon Europe): \textbf{React}-Web-Frontend über C++-Inferenz auf GPU-Servern und einer \textbf{Berghof-SPS mit Hailo}
            % Vor "25 Partner" (2026-09-22), zum Nachschlagen erhalten:
            % \item \textbf{Aktivitätserkennung} für \link{https://remanet-project.eu}{RemaNet} (Horizon Europe): Kamera-Pipelines mit \textbf{React}-Web-Frontend über C++-Inferenz auf GPU-Servern
            % Mit "25 Partner", vor der Umformulierung des Inhabers (2026-09-22),
            % zum Nachschlagen erhalten:
            % \item \textbf{Aktivitätserkennung} für \link{https://remanet-project.eu}{RemaNet} (Horizon Europe, 25 Partner): Kamera-Pipelines mit \textbf{React}-Web-Frontend über C++-Inferenz auf GPU-Servern
            % Siehe die Anmerkung im englischen Zweig: ONNX-Inferenz und
            % PyTorch-Training auf NVIDIA-GPUs.
            \def\cvFhgRemanet{\ifcvshowremanet
            \item \textbf{Aktivitätserkennung} für \link{https://remanet-project.eu}{RemaNet} (Horizon Europe, 25 Partner): Kamera-Pipelines mit \textbf{React}-Web-Frontend, \textbf{ONNX}-Inferenz\slash\textbf{PyTorch}-Training auf NVIDIA-GPUs
            \fi}
            % Vor Bildrate und Framework (2026-09-22), zum Nachschlagen erhalten:
            % \item \textbf{Objekterkennung} auf einer eingebetteten \textbf{Berghof-SPS}, die ich mit einem \textbf{Hailo}-KI-Beschleuniger nachgerüstet habe
            \def\cvFhgBerghof{\ifcvshowberghof
            \item In PyTorch trainierte \textbf{Objekterkennung} mit \textbf{$\sim$14 fps} auf einer eingebetteten \textbf{Berghof-SPS}, die ich mit einem \textbf{Hailo}-KI-Beschleuniger nachgerüstet habe\ifcvshowmldetail; \textbf{INT8}-quantisiert und mit Industriekunden anhand von Feldfeedback iteriert\fi
            \fi}
            % Frühere Fassung, 2026-09-22, zum Nachschlagen erhalten:
            % \item \textbf{Personenerkennung} für industrielle Sicherheit: \textbf{4 Kameras} parallel, Inferenzen \textbf{gebatcht}
            \def\cvFhgSafety{\ifcvshowsafety
            \item \textbf{Personenerkennung} für die Sicherheitsanwendung eines Industriekunden geliefert: \textbf{4 Kameras} parallel, Inferenzen \textbf{gebatcht}
            \fi}
            % Siehe die Anmerkung im englischen Zweig: KI-PrOT, Ziel und nicht
            % Messergebnis, daher "um ... zu sparen".
            \def\cvFhgKiprot{\ifcvshowkiprot
            \item \textbf{Visuelle Klassifikation} von Bauteilen in Verschleppungsklassen für \link{https://ki-prot.de}{KI-PrOT}, um Chemikalien in der Oberflächentechnik zu sparen: Datensatz erstellt, annotiert und mit \textbf{DVC} versioniert
            \fi}
            % Bis hier ist alles nur definiert; diese Zeile gibt es aus, in der
            % Reihenfolge aus cv.tex oder dem Profil.
            \cvFraunhoferOrder
        \end{itemize}
    }
    % Siehe die Anmerkung im englischen Zweig (PyTorch und uv seit
    % 2026-09-22). Frühere Listen, zum Nachschlagen erhalten:
    % {Python, Go, C++, CUDA, CMake, Bazel, Meson, vLLM, LoRA, GitLab CI/CD}
    % {Python, Go, C++, CUDA, PyTorch, CMake, Bazel, Meson, vLLM, LoRA, GitLab CI/CD}
    % {Python, Astral uv, Go, C++, CUDA, PyTorch, CMake, Bazel, Meson, vLLM, LoRA, GitLab CI/CD}
    {Python, uv, Go, C++, CUDA, PyTorch, CMake, Bazel, Meson, vLLM, LoRA, GitLab CI/CD}
        
    \emptySeparator
    
    \experience
    {\DTMlangsetup*{showdayofmonth=false}\DTMdisplaydate{2023}{05}{01}}   {Research assistant}{Karlsruhe}{CAS Software AG}
    {\DTMlangsetup*{showdayofmonth=false}
    \DTMdisplaydate{2022}{10}{15}} 
    {
        \ifcvshowcasdetail
        \begin{itemize}
            % Siehe die Anmerkung im englischen Zweig.
            % Nominal statt verbal, weil kürzer: die verbale Fassung ("Training
            % und Deployment von ONNX-Modellen für ...") lief um vier Zeichen
            % über die Zeile und kostete eine zweite für das Wortende.
            % Siehe die Anmerkung im englischen Zweig.
            \item ONNX-Modelle und Eye Tracking für Echtzeit-Benutzerführung in Mixed Reality
        \end{itemize}
        \fi
    }
    {Machine learning, ONNX, Mixed reality, HoloLens2, C\#, Python}
    %\emptySeparator
    \end{experiences}
}
%
    % Four years of KIT tutoring (2018-2022). Back in after the two GLM-5.2
    % bullets were merged: the duplication they carried paid for part of the
    % ~5 lines this costs. Worth it for two reasons -- it is the only entry
    % covering 2018-2022, which is what makes a "4+ years" requirement legible,
    % and the posting asks for a "collaborative, low-ego" engineer, which
    % teaching is the evidence for. The heading stays off (see cv.tex), so it
    % reads as one more job under Experiences, which is what it is.
    % YAAC Another Awesome CV LaTeX Template
%
% This template has been downloaded from:
% https://github.com/darwiin/yaac-another-awesome-cv
%
% Author:
% Christophe Roger
%
% Template license:
% CC BY-SA 4.0 (https://creativecommons.org/licenses/by-sa/4.0/)

%Section: Scholarships and additional info
\IfLanguageName{english}{

    % Heading is optional: a profile that inputs this file straight after
    % section_experience wants the entry to read as one more job under
    % Experiences (which it is -- a paid research-assistant post), not as a
    % section of its own. See \ifcvshowteachingtitle in cv.tex.
    \ifcvshowteachingtitle
    \sectionTitle{Teaching and Mentoring}{\faUsers}
    \fi
    
    \begin{experiences}

        \experience
        {\DTMlangsetup*{showdayofmonth=false}
        % Third argument is the location slot (the class sets it in small caps
        % beside the institution) -- it held "Research Assistant", which put the
        % employment type where every other entry has a city and wrapped the
        % header onto a second line. The role is already in the title.
        \DTMdisplaydate{2022}{08}{01}}   {Tutor | Technical Computer Science}{Karlsruhe}{Karlsruhe Institute of Technology}
        {\DTMlangsetup*{showdayofmonth=false}
        \DTMdisplaydate{2018}{10}{01}} {
        \begin{itemize}
            % One bullet, not the full curriculum: this section only earns its
            % place on a one-page CV if it costs ~4 lines, and outside academia
            % the syllabus detail is not what the section is evidence for --
            % four years of explaining hard material to people who do not have
            % the context yet is. The original enumeration is kept below;
            % restore it for teaching-oriented applications.
            % The syllabus enumeration goes when the section has to pay for
            % itself, which the note above already anticipated: the subject and
            % the four years are what the entry is evidence for, not the topic
            % list. The full list survives in the commented lines below.
            % "while studying" is the whole point of the line: four years of
            % teaching *alongside his own degree* is an energy fact, and it is
            % free -- the entry title already says Tutor at KIT, so dropping
            % "undergraduate" costs nothing and buys the two words.
            \item Taught digital technology and computer organization for four years while studying
            % \item Teaching undergraduates in digital technology. This includes inter alias control units, run-time effects, computer arithmetic, boolean algebra, NAN/NOR gates, MOSFET, CMOS, Quine-McClusky, hazards,
            % sequential circuit
            % \item Teaching undergraduates in computer organization which includes inter alias c programming, caches, assembler, memory organization, DLX-pipelining, RAM, paging, segmentation
        \end{itemize}
        }
        {\LaTeX, C, Assembler, MIPS, RISC-V}
        \emptySeparator
        % Commented out to fit the CV on one page -- restore by uncommenting.
        %\experience
        %{\DTMlangsetup*{showdayofmonth=false}
        %\DTMdisplaydate{2021}{08}{01}}   {Mathematics | Tutor}{Freelancer}{Karlsruhe}
        %{\DTMlangsetup*{showdayofmonth=false}
        %\DTMdisplaydate{2020}{01}{01}} { %August 2021
        %\begin{itemize}
        %    \item Teaching pupils in mathematics
        %\end{itemize}
        %}
        %{Mathematics}
        %\emptySeparator
        %\experience{\DTMlangsetup*{showdayofmonth=false}
        %\DTMdisplaydate{2015}{8}{31}}
        %{Federal volunteer service}
        %{Supervisor for a person with autism}
        %{Schwäbisch Hall}
        %{\DTMlangsetup*{showdayofmonth=false}
        %\DTMdisplaydate{2015}{02}{16}}
        %{
        %\begin{itemize}
        %  \item Assisting a 16 year old boy with autism in school.
        %\end{itemize}
        %}
        %{Volunteering}
    \end{experiences}        
}{
    % Siehe die Anmerkung im englischen Zweig.
    \ifcvshowteachingtitle
    \sectionTitle{Lehre und Mentoring}{\faUsers}
    \fi
    
    \begin{experiences}
        \experience
        {\DTMlangsetup*{showdayofmonth=false}
        % Siehe die Anmerkung im englischen Zweig zum Orts-Argument.
        \DTMdisplaydate{2022}{08}{01}}   {Tutor | Technische Informatik}{Karlsruhe}{Karlsruher Institut für Technologie}
        {\DTMlangsetup*{showdayofmonth=false}
        \DTMdisplaydate{2018}{10}{01}} {
        \begin{itemize}
            % Siehe die Anmerkung im englischen Zweig: eine Zeile statt des
            % vollen Curriculums, das unten erhalten bleibt.
            % "neben dem Studium" gestrichen: der Zeitraum steht links daneben
            % und überschneidet sich sichtbar mit dem Studium. Die Wörter waren
            % genau die, die "Speicherverwaltung" auf eine dritte Zeile schoben.
            % Siehe die Anmerkung im englischen Zweig.
            % Siehe die Anmerkung im englischen Zweig.
            \item Vier Jahre Lehre in Technischer Informatik und Rechnerorganisation neben dem Studium
            % \item Unterrichten von Studierenden in Grundlagen der Mikroelektronik, Entwurf und Aufbau von einfachen datenverarbeitenden Systemen, logischen Schaltnetzen und Schaltwerken
            % \item Weitere Themen: boolesche Algebra, boolesche Funktionen, Minimierungsverfahren, Entwurf einfacher arithmetischer Einheiten, Grundlagen des Aufbaus und Organisation von Rechnern, Assemblerprogrammierung, Pipelining, Cache-Speicher, virtuelle Speicherverwaltung, etc.
        \end{itemize}
        }
        {\LaTeX, C, Assembler, MIPS, RISC-V}
        \emptySeparator
        % Commented out to fit the CV on one page -- restore by uncommenting.
        %\experience
        %{\DTMlangsetup*{showdayofmonth=false}
        %\DTMdisplaydate{2021}{08}{01}}   {Mathematik | Tutor}{Selbstständig}{Karlsruhe}
        %{\DTMlangsetup*{showdayofmonth=false}
        %\DTMdisplaydate{2020}{01}{01}} { %August 2021
        %\begin{itemize}
        %    \item Nachhilfe für Schüler in Mathematik
        %\end{itemize}
        %}
        %{Mathematik}
        %\emptySeparator
        %\experience{\DTMlangsetup*{showdayofmonth=false}
        %\DTMdisplaydate{2015}{8}{31}}
        %{Bundesfreiwilligendienst}
        %{Schulbegleitung für eine Person mit Autismus}
        %{Schwäbisch Hall}
        %{\DTMlangsetup*{showdayofmonth=false}
        %\DTMdisplaydate{2015}{02}{16}}
        %{
        %\begin{itemize}
        %  \item Ünterstützung eines 16-Jährigen in der Schule
        %\end{itemize}
        %}
        %{Volunteering}
        
    \end{experiences}
}
  
%
    % YAAC Another Awesome CV LaTeX Template
%
% This template has been downloaded from:
% https://github.com/darwiin/yaac-another-awesome-cv
%
% Author:
% Christophe Roger
%
% Template license:
% CC BY-SA 4.0 (https://creativecommons.org/licenses/by-sa/4.0/)

%Section: Scholarships and additional info
\IfLanguageName{english}{
    
    \sectionTitle{Education}{\faMortarBoard}
    \begin{experiences}
    \education
    {\DTMlangsetup*{showdayofmonth=false}
    \DTMdisplaydate{2023}{09}{30}}   {Master of Science Computer Science}{Karlsruhe}{Karlsruhe Institute of Technology}
    {\DTMlangsetup*{showdayofmonth=false}
    \DTMdisplaydate{2020}{04}{01}}
    {Overall grade: 1.9/ US-GPA: 3.1}
    {
        \begin{itemize}
            \item Thesis: Designing User-adaptive Guidance for Mixed Reality Using Eye Tracking
        \end{itemize}
    }
    {AI, Computer Graphics and Geometry Processing, Anthropomatics and Cognitive Systems}
    \emptySeparator
    \education{\DTMlangsetup*{showdayofmonth=false}
    \DTMdisplaydate{2020}{03}{01}}
    {Bachelor of Science Computer Science}
    {Karlsruhe}
    {Karlsruhe Institute of Technology}
    { \DTMlangsetup*{showdayofmonth=false}
    \DTMdisplaydate{2015}{10}{01}}
    {Overall grade: 2.5/ US-GPA: 2.5}
    {
    \begin{itemize}
        \item Thesis: Temporally Stable Blue Noise Error Distribution in Screen Space for Real Time Applications
    \end{itemize}
    }
    {Path Tracing, Blue Noise, C++, Real-Time Systems, Low Latency, High Performance}
    % A-levels entry commented out to pay for the Open Source section --
    % restore by uncommenting, and drop something else instead.
%     \emptySeparator
%     \education{2014}
%     {Pupil | Higher education entrance qualification (A-levels)}
%     {Schwäbisch Hall}
%     {Gymnasium bei St. Michael}
%     {2006}
%     {Overall grade: 2.2/ US-GPA: 2.8}
%     {
%         % Commented out to fit the CV on one page -- restore by uncommenting.
%         %\begin{itemize}
%         %    \item First steps in Procedural programming: Sudoku project
%         %    \item Started learning Java autodidactically
%         %\end{itemize}
%     }
%     {Delphi, Pascal, Java}
    \end{experiences}
    
}{
    
    \sectionTitle{Bildung}{\faMortarBoard}
    \begin{experiences}
    % german
    \education
    {\DTMlangsetup*{showdayofmonth=false}
    \DTMdisplaydate{2023}{09}{30}}   {Master of Science Informatik}{Karlsruhe}{Karlsruher Institut für Technologie}
    {\DTMlangsetup*{showdayofmonth=false}
    \DTMdisplaydate{2020}{04}{01}}
    {Gesamtnote: 1,9}
    {
        \begin{itemize}
            \item Thesis: Entwurf von benutzeranpassbarem Szeneninhalt für Mixed Reality mit Hilfe von Augen- und Handnachverfolgung
        \end{itemize}
    }
    {AI, Computergrafik und Geometrieverarbeitung, Anthropomatik und Kognitive Systeme}
    \emptySeparator
    \education{\DTMlangsetup*{showdayofmonth=false}
    \DTMdisplaydate{2020}{03}{01}}
    {Bachelor of Science Informatik}
    {Karlsruhe}
    {Karlsruher Institut für Technologie}
    { \DTMlangsetup*{showdayofmonth=false}
    \DTMdisplaydate{2015}{10}{01}}
    {Gesamtnote: 2,5}
    {
    \begin{itemize}
        \item Thesis: Zeitlich stabile Blue Noise Fehlerverteilung im Bildraum für Echtzeitanwendungen
    \end{itemize}
    }
    {Path Tracing, Blue Noise, C++, Echtzeitsysteme, Low Latency, High Performance}
    % A-levels entry commented out to pay for the Open Source section --
    % restore by uncommenting, and drop something else instead.
%     \emptySeparator
%     \education{2014}
%     {Schüler | Allgemeine Hochschulreife}
%     {Schwäbisch Hall}
%     {Gymnasium bei St. Michael}
%     {2006}
%     {Gesamtnote: 2,2}
%     {
%         % Commented out to fit the CV on one page -- restore by uncommenting.
%         %\begin{itemize}
%         %    \item erste Schritte in proceduraler Programmierung: Sudokuprojekt
%         %    \item Java autodidaktisch gelernt
%         %\end{itemize}
%     }
%     {Delphi, Pascal, Java}
    
    \end{experiences}
}
  


%
    % Languages (French B1 -- the role is Paris-based) and the ECML PKDD
    % publication. Cheap in vertical space: the row's height is set by its
    % tallest column, and Languages sets it either way.
    % YAAC Another Awesome CV LaTeX Template
%
% This template has been downloaded from:
% https://github.com/darwiin/yaac-another-awesome-cv
%
% Author:
% Christophe Roger
%
% Template license:
% CC BY-SA 4.0 (https://creativecommons.org/licenses/by-sa/4.0/)

% Languages / Hobbies / Publications share one \threecolumnsection row.
% \small and the dropped \vspace{1em} under each title are what keep that row
% inside the space left on page one -- the CV is a strict one-pager.
\IfLanguageName{english}{
    %Section: Languages
    \par\noindent
    \small
    \threecolumnsection
    {\sectionTitle{Languages}{\faLanguage}
    \begin{itemize}
    	\item German (native)
    	\item English, level C1
    	\item French, level B1
    	\item Latin
    \end{itemize}}
    {\sectionTitle{Hobbies}{\faBicycle}
    \begin{itemize}
        \item Powerlifting
        % Commented out per request -- restore by uncommenting.
        % \item Bouldering
        \item Guitar
        \item Philosophy
        \item Politics
    \end{itemize}
    }
    {\sectionTitle{Publications}{\faFileTextO}
    % \raggedright as in \twocolumnsection: justifying one long title across
    % this wide column stretches the interword spaces visibly.
    \raggedright
    \begin{itemize}
        \item \textquote{Semantic Robustness Probing via Inpainting} -- ECML PKDD 2026 Demo Track (with N.\ Steckhan)
    \end{itemize}
    }
}{
    %Section: Languages
    \par\noindent
    \small
    \threecolumnsection
    {\sectionTitle{Sprachen}{\faLanguage}
    \begin{itemize}
    	\item Deutsch (Muttersprache)
    	% CEFR levels in parentheses: ", Level B1" overflows this narrow column.
    	\item Englisch (C1)
    	\item Französisch (B1)
    	\item Latein
    \end{itemize}}
    {\sectionTitle{Hobbies}{\faBicycle}
    \begin{itemize}
        \item Powerbuilding
        \item Gitarre
        \item Philosophie
        \item Politik
    \end{itemize}
    }
    {\sectionTitle{Publikationen}{\faFileTextO}
    % \raggedright as in \twocolumnsection: justifying one long title across
    % this wide column stretches the interword spaces visibly.
    \raggedright
    \begin{itemize}
        \item \textquote{Semantic Robustness Probing via Inpainting} -- ECML PKDD 2026 Demo Track (mit N.\ Steckhan)
    \end{itemize}
    }
}
%
    % Left out on purpose, not for lack of relevance -- one page. The GitHub
    % link in the header and in the summary carries the project evidence:
    % \input{section_projects}%
    % \IfLanguageName{english}{
    \sectionTitle{Technical Skills}{\faCogs}
    \twocolumnsection
    {\textbf{Programming \& Markup}
    \vspace{0.5em}
    \begin{itemize}[leftmargin=*]
        \item C, C++
        \item Python
        \item Rust
        \item GLSL, HLSL
        \item C\#, Dart
        \item LaTeX, Markdown
    \end{itemize}}
    % Commented out to fit the CV on one page -- restore by uncommenting
    % (and switching back to \threecolumnsection with this as the 2nd arg).
    % {\textbf{Graphics \& Rendering}
    % \vspace{0.5em}
    % \begin{itemize}[leftmargin=*]
    %     \item Vulkan, OpenGL
    %     \item PBR, Ray Tracing
    %     \item Deferred Shading
    %     \item Shadow Mapping
    %     \item SSAO, TAA, SSR
    %     \item Compute Shaders
    % \end{itemize}}
    {\textbf{ML \& Tools}
    \vspace{0.5em}
    \begin{itemize}[leftmargin=*]
        \item PyTorch, ONNX
        \item Mixed Reality (HoloLens2)
        \item CMake, Docker
        \item CI/CD, Cross-Compilation
        \item Git, Linux
        \item Unity, Blender
    \end{itemize}}
}{
    \sectionTitle{Technische Fähigkeiten}{\faCogs}
    \twocolumnsection
    {\textbf{Programmiersprachen}
    \vspace{0.5em}
    \begin{itemize}[leftmargin=*]
        \item C, C++
        \item Python
        \item Rust
        \item GLSL, HLSL
        \item C\#, Dart
        \item LaTeX, Markdown
    \end{itemize}}
    % Commented out to fit the CV on one page -- restore by uncommenting
    % (and switching back to \threecolumnsection with this as the 2nd arg).
    % {\textbf{Grafik \& Rendering}
    % \vspace{0.5em}
    % \begin{itemize}[leftmargin=*]
    %     \item Vulkan, OpenGL
    %     \item PBR, Ray Tracing
    %     \item Deferred Shading
    %     \item Shadow Mapping
    %     \item SSAO, TAA, SSR
    %     \item Compute Shader
    % \end{itemize}}
    {\textbf{ML \& Tools}
    \vspace{0.5em}
    \begin{itemize}[leftmargin=*]
        \item PyTorch, ONNX
        \item Mixed Reality (HoloLens2)
        \item CMake, Docker
        \item CI/CD, Cross-Compilation
        \item Git, Linux
        \item Unity, Blender
    \end{itemize}}
}
%
}
